\section{Problem and Adaptive Opportunity}
\label{sec:problem}

\paragraph{Control and work.}
A target run consists of $U_\star$ prescribed optimizer updates, executed as
sequential checkpointed chunks. At boundary $k$, time is $t_k$ and remaining
work is $U_k$. The action is a feasible node count
$n_k\in\mathcal A=\{4,16,64,128\}$. The next chunk is submitted immediately;
there is no intentional deferral, concurrent reservation, or resizing of a
running chunk. Scale choices may rise or fall between allocations; ``down
means using fewer than the maximum supported nodes, not irreversible shrinking. We use $H=48$ hours as the study's maximum requested wall time,
subject to the chosen partition's limits; it is not a universal Slurm limit.

For sharded or replicated data parallelism without tensor or pipeline
parallelism, let $g$ be GPUs per node ($g=8$ for the published AMSP profiles),
$b$ the per-GPU micro-batch, and $a_n$ accumulation. Each
action satisfies
\begin{equation}
 B_\star=ngba_n,\qquad a_n=B_\star/(ngb)\in\mathbb N.
 \label{eq:batch}
\end{equation}
The same data order, sequence length, optimizer-update target, and update-indexed
learning-rate schedule apply across actions. Accumulation implements the fixed
batch~\cite{shoeybi2019megatron}; it does not guarantee bitwise equality or equal
final quality. A configuration that fails the batch, memory, or partition
constraints is masked for every method.

Published steady-state throughput determines $q_n$, the updates per hour.
Initialization/restart $r_n$ and complete-checkpoint time $h_n$ are declared
scenarios, not measurements supplied by the throughput source. Define
\begin{align}
 v_{k,n}&=\min\{U_k,\lfloor q_n(H-r_n-h_n)\rfloor\},\nonumber\\
 d_{k,n}&=r_n+v_{k,n}/q_n+h_n,\label{eq:chunk}\\
 U_{k+1}&=U_k-v_{k,n_k},\quad
 t_{k+1}=t_k+w_k+d_{k,n_k}.\nonumber
\end{align}
Only positive-progress actions are allowed. The request's wall time is
$d_{k,n}$ rounded up to the scheduler's time granularity, capped at $H$; the
job releases resources on completion. Thus the final partial chunk has a
shorter request, which can change admission. The assumed restart contract includes model,
optimizer, learning-rate state, global update index, random-state metadata,
and the next global sample position; this study does not implement or validate
AMSP state migration. The final output save is included in
$h_n$. End-to-end turnaround time is $T=t_K-t_0$, including every queue wait
and allocation interval.

\paragraph{A fixed-scale null case.}
The following observation defines what scaling profiles alone can establish.

\begin{proposition}[No improvement below the fixed hull]
\label{prop:hull}
Suppose update rates and per-node powers are constant for each scale, carbon
intensity is constant, and there are no waits or restart/checkpoint overheads.
Every execution of the same work, including one that switches scales, has a
time--carbon pair in the convex hull of the fixed-scale pairs.
\end{proposition}
\noindent\emph{Proof.}
Let $x_n$ be the fraction of updates executed at scale $n$, so
$x_n\geq0$ and $\sum_nx_n=1$. If $T_n=U_\star/q_n$ and $C_n$ are the
fixed-scale totals, additivity gives
$(T,C)=\sum_nx_n(T_n,C_n)$.\hfill$\square$

This is a statement about mean time--carbon geometry, not a per-job deadline
guarantee for randomized fixed policies. A preplanned sequence can complete
individual runs at an intermediate time--carbon point even when a per-run
fixed mixture violates a deadline. Thus beating a feasible fixed mixture
alone does not prove that queue feedback is useful. Under real queueing, a request changes
both its current wait and the time and scheduler state at which the next
request arrives; the linear decomposition need no longer hold. Likewise,
time-varying carbon intensity makes an update's carbon depend on its execution
window. These effects create an opportunity, not a guarantee of an adaptive
gain. Our main comparison retains all fixed scales, a deployable mixture,
and a plan-once controller; action-switch frequency alone is not evidence
of improvement.

\paragraph{What the slider can establish.}
Let $V_{\cal P}(D)$ be the infimum of the endpoint objective in
Equation~\ref{eq:objective} over feasible policies in class $\cal P$. These classes describe admissible
controllers with history, not the expressiveness of a particular network.
Fixed policies and their per-run mixtures are contained in preplanned
sequences, which are contained in feedback policies. Therefore
$V_{\rm feedback}(D)\leq V_{\rm preplanned}(D)\leq V_{\rm fixed}(D)$.
Each ideal envelope is nonincreasing in $D$, since relaxing the budget retains
every formerly feasible policy. This is an opportunity bound, not a guarantee
that finite PPO training finds the envelope. Strict improvement requires
useful timing or admission variation; the null case above rules out claiming
it from scaling inefficiency alone. A sequence that beats both fixed points
still needs a preplanned comparator to establish the value of feedback.
